\documentclass[notes=show]{beamer}
\usepackage{mathpazo}
\usepackage{hyperref}
\usepackage{multimedia}
\usepackage{tcolorbox}
\usepackage{tikz}
\usetikzlibrary{shadows}
\usetheme{metropolis}
\usecolortheme{default}
\setbeamertemplate{itemize item}{%
    \begin{tikzpicture}
        \shade[ball color=black!100!yellow, preaction={fill=black,
opacity=.00,transform canvas={xshift=1mm,yshift=-1mm, yscale=0.5}}] (0,0) circle (0.6ex);
    \end{tikzpicture}
}
\setbeamertemplate{itemize subitem}{%
    \begin{tikzpicture}
        \shade[ball color=black!100!white, preaction={fill=black,
        opacity=.00,transform canvas={xshift=1mm,yshift=-1mm, yscale=0.5}}] (0,0) circle (0.6ex);
    \end{tikzpicture}
}
\setbeamercolor{frametitle}{bg=white}
\setbeamercolor{frametitle}{fg=black}
\setbeamercolor{background canvas}{bg=white}
\setbeamercolor{block body}{bg=mDarkTeal!30}
\setbeamercolor{block title}{bg=mDarkTeal,fg=black!2}
\setbeamertemplate{navigation symbols}{}
\setbeamertemplate{footline}[page number]
\newenvironment{stepenumerate}{\begin{enumerate}[<+->]}{\end{enumerate}}
\newenvironment{stepitemize}{\begin{itemize}[<+->]}{\end{itemize} }
\newenvironment{stepenumeratewithalert}{\begin{enumerate}[<+-| alert@+>]}{\end{enumerate}}
\newenvironment{stepitemizewithalert}{\begin{itemize}[<+-| alert@+>]}{\end{itemize} }

\graphicspath{{figures/}}

\begin{document}

\title{Firms and Labor Market Inequality: Evidence and Some Theory}
\subtitle{}
\date{Card, Cardoso, Heining, Kline \bigskip \\
\textit{Journal of Labor Economics}, 2018, 36:S1, S13-S70}
\author{}
\maketitle

%%%%%%%%%%%%%%%%%%%%%%%%%%%%%%%%%%%%%%%%%%%%%%%%%%%%%%%%%
\section{Introduction}
%%%%%%%%%%%%%%%%%%%%%%%%%%%%%%%%%%%%%%%%%%%%%%%%%%%%%%%%%

\begin{frame}{Inequality in imperfectly competitive labor markets}
\begin{itemize}
\item Employers set wages (Robinson 33, Slichter 50, Krueger\&Summers 88) \medskip
\item Difficult to distinguish firm-wage premia from sorting (Gibbons\&Katz 92) \medskip
\item Impact of shocks to firm productivity on wages of stayers finding rent-sharing elasticities of 0.05-0.15 \medskip
\item AKM using switchers to identify additive worker and firm FE finding that 20\% of variances in wages is due to firm FE
\end{itemize}
\end{frame}

\begin{frame}{Linking rent-sharing elasticities to firm FE and worker sorting}
\begin{itemize}
\item Review of literature on rent-sharing and studies using AKM \medskip
\item Use data for Portugal to show that more productive firms have higher firm FE \medskip
\item Model in which heterogeneous firms set wages when workers value firms differently \medskip 
\item Model has become known as ``new classical monopsony'' \medskip
\item Show that this model explains rent-sharing elasticities, firm FE and sorting of workers across firms.
\end{itemize}
\end{frame}

%%%%%%%%%%%%%%%%%%%%%%%%%%%%%%%%%%%%%%%%%%%%%%%%%%%%%%%%%
\section{Productivity, Wages, and Rent Sharing}
%%%%%%%%%%%%%%%%%%%%%%%%%%%%%%%%%%%%%%%%%%%%%%%%%%%%%%%%%

\newpage
\begin{center}
\begin{figure}
\includegraphics[height=8.5cm,width=\textwidth]{figures/fig1.pdf}
\end{figure} 
\end{center}
\newpage

\begin{frame}{Impartial pass-through of productivity to wages}
\begin{itemize}
\item Literature on rent sharing examines whether productivity differences feed into wage differences \medskip
\item Rent-sharing elasticities of 0.05-0.15 show impartial pass-through suggesting firms have wage setting power \medskip
\item Analog of literature in international economics and IO finding impartial pass-through of cost shocks to prices suggesting firms have price setting power \medskip
\item Rent-sharing elasticities are sensitive to measurement issues
\end{itemize}
\end{frame}

%%%%%%%%%%%%%%%%%%%%%%%%%%%%%%
\subsection*{Measuring Rents}
%%%%%%%%%%%%%%%%%%%%%%%%%%%%%%

\begin{frame}{Measuring rents}
\begin{itemize}
\item Assume that the wage at firm $j$ is given by:
\begin{equation*}
    w_{j} = b + \gamma \frac{Q_{j}}{N_{j}} \tag{1}
\end{equation*}
with $Q_{j}/N_{j}$ quasi-rents per worker. 
\medskip
\item Quasi-rents are given by:
\begin{equation*}
    Q_{j} = VA_{j} - bN_{j} - rK_{j}
\end{equation*}
with $VA_{j}$ value added (revenue net of material costs). \medskip
\item Value added is given by:
\begin{equation*}
    VA_{j} = P_{j} T_{j} f(N_{j},K_{j})
\end{equation*}
with $T_{j}$ firm-specific $TFP_{j}$.
\end{itemize}
\end{frame}

\begin{frame}{Rent-sharing elasticities with quasi-rents and profit per worker}
\begin{itemize}
\item The elasticity of wages w.r.t. quasi-rents per worker is:
\begin{equation*}
    \xi_{Q{j}} \equiv \frac{d w_{j}}{d (Q_{j}/N_{j})} \frac{Q_{j}/N_{j}}{w_{j}} = \frac{\gamma [Q_{j}/N_{j}]}{b + \gamma [Q_{j}/N_{j}]} \tag{2}
\end{equation*}
which corresponds to the share of quasi-rents in wages. \medskip
\item $\xi_{Q{j}} $ is also elasticity of wages w.r.t. profit per worker because:
\begin{align*}
    \frac{\Pi_{j}}{N_{j}} & = \frac{1}{N_{j}} \left[ VA_{j} - w_{j}N_{j} - rK_{j} \right] \nonumber \\
    & = \frac{1}{N_{j}} \left[ VA_{j} - bN_{j} - rK_{j} - (w_{j} - b)N_{j} \right] \nonumber \\
    & = \frac{1}{N_{j}} \left[ Q_{j} - \gamma \frac{Q_{j}}{N_{j}}N_{j} \right] \nonumber \\
    & = \frac{Q_{j}}{N_{j}} \left[ 1 - \gamma \right] \nonumber
\end{align*}
\end{itemize}
\end{frame}

\begin{frame}{Rent-sharing elasticities with value added per worker}
\begin{itemize}
\item The elasticity of wages w.r.t. value added per worker is:
\begin{align*}
    \xi_{j} & \equiv \frac{d w_{j}}{d (VA_{j}/N_{j})} \frac{VA_{j}/N_{j}}{w_{j}} \nonumber \\
    & = \frac{d w_{j}}{d (Q_{j}/N_{j})} \frac{VA_{j}/N_{j}}{w_{j}} \nonumber \\
    & = \frac{d w_{j}}{d (Q_{j}/N_{j})} \frac{Q_{j}/N_{j}}{w_{j}} \frac{VA_{j}}{Q_{j}} \nonumber \\
    & = \xi_{Q{j}} \frac{VA_{j}}{Q_{j}} \nonumber
\end{align*}
\item $\xi_{j} >  \xi_{Q{j}} $ because $ VA_{j} > Q_{j}$ (e.g. $ VA_{j} / Q_{j} = 2$) \medskip 
\item A majority of studies estimates $\xi_{j} $
\end{itemize}
\end{frame}

%%%%%%%%%%%%%%%%%%%%%%%%%%%%%%
\subsection*{A Summary of the Rent-Sharing Literature}
%%%%%%%%%%%%%%%%%%%%%%%%%%%%%%

\newpage
\begin{frame}{Summary of the rent-sharing literature}
\begin{center}
\begin{figure}
\includegraphics[scale=0.98]{Cardetal/figures/tab1-1.pdf}
\end{figure} 
\end{center}
\end{frame}

\newpage
\begin{frame}{Summary of the rent-sharing literature}
\begin{center}
\begin{figure}
\includegraphics[scale=0.98]{Cardetal/figures/tab1-2.pdf}
\end{figure} 
\end{center}
\end{frame}

%%%%%%%%%%%%%%%%%%%%%%%%%%%%%%
\subsection*{Specification Issues: A Replication in Portuguese Data}
%%%%%%%%%%%%%%%%%%%%%%%%%%%%%%

\begin{frame}{Replication using Portuguese data}
\begin{itemize}
\item Wage data from Quadros de Pessoal (QP), a census of private sector employees between 2005-2009. \medskip 
\item Firm-specific financial information from SABI database by Bureau van Dijk with information on value added and sales per worker between 2004-2010. \medskip
\item Merge QP and SABI using information on detailed location, industry, firm creation date, shareholder equity, and annual sales that are available in both data sets. \medskip
\item Regress wages onto value added or sales per worker controlling for human capital variables, industry and location.
\end{itemize}
\end{frame}

\newpage
\begin{frame}{Rent-sharing elasticities for Portuguese male workers}
\begin{center}
\begin{figure}
\includegraphics[scale=0.95]{Cardetal/figures/tab2-1.pdf}
\end{figure} 
\end{center}
\end{frame}

\newpage
\begin{frame}{Rent-sharing elasticities for Portuguese male workers}
\begin{center}
\begin{figure}
\includegraphics[scale=0.95]{Cardetal/figures/tab2-2.pdf}
\end{figure} 
\end{center}
\end{frame}

%%%%%%%%%%%%%%%%%%%%%%%%%%%%%%%%%%%%%%%%%%%%%%%%%%%%%%%%%
\section{Firm Switching}
%%%%%%%%%%%%%%%%%%%%%%%%%%%%%%%%%%%%%%%%%%%%%%%%%%%%%%%%%

\begin{frame}{AKM model}
%Data set contains N* person-year observations on N workers and J establishments \\
\begin{itemize} 
    \item Log daily real wage $y_{it}$ of individual $i$ in year $t$ is given by:
    \begin{equation*}
	\ln(w_{it}) = \alpha_i + \psi_{J(i,t)} + x'_{it}\beta + \epsilon_{it}
    \end{equation*}
\item Estimates of firm effects typically explain 15\%–25\% of the variance of wages \medskip
\item Sorting of workers and firms with high FE also explains part of the variance in wages \medskip
\item Problem is that the person and firm effects are overstated and sorting is understated in small samples and limited mobility
\end{itemize}
\end{frame}

\begin{frame}{Variance decompositions}
\begin{itemize}
\item Variance of wages can be decomposed as:
    \begin{align*}
	    Var(\ln(w_{it}))& = Var(\alpha_i)+Var(\psi_{J(i,t)})+Var(x_{it}'\beta) \tag{3} \\
				&+ 2Cov(\alpha_i,\psi_{J(i,t)})+2Cov(\psi_{J(i,t)},x_{it}'\beta)  \\
				&+ 2Cov(\alpha_i,x_{it}'\beta) + Var(\epsilon_{it})
	\end{align*}
\item An alternative decomposition is:
    \begin{align*}
		Var(\ln(w_{it}))& = Cov(\ln(w_{it},\alpha_i) + Cov(\ln(w_{it},\psi_{J(i,t)}) \tag{4} \\
				& + Cov(\ln(w_{it},x_{it}'\beta) + Cov(\ln(w_{it}, \epsilon_{it})
    \end{align*}
which is an ensemble assessment of each component, e.g.:
\begin{align*}
    Cov(\ln(w_{it}),\psi_{J(i,t)}) = & Var(\psi_{J(i,t)}) + Cov(\alpha_{i},\psi_{J(i,t)}) \\
    & + Cov(x_{it}'\beta,\psi_{J(i,t)})
\end{align*}
\end{itemize}
\end{frame}

%%%%%%%%%%%%%%%%%%%%%%%%%%%%%%
\subsection*{Identifying Age and Time Effects}
%%%%%%%%%%%%%%%%%%%%%%%%%%%%%%

\begin{frame}{Identifying age and time effects}
\begin{itemize}
\item A technical issue that arises in AKM models is the linear inclusion of age $a_{it}$ in $X_{it}$. \medskip
\item It is standard to include year effects $t$ in $X_{it}$. \medskip
\item This raises an identification problem because $a_{it} = t - b_{i}$ with $b_{i}$ birth year which loads onto $\alpha_{i}$ \medskip
\item Use actual experience instead of age if available and if gaps are exogenous conditional on $\alpha_{i}$. \medskip
\item Impose restrictions on age or time effects. 
\end{itemize}
\end{frame}

\newpage
\begin{frame}{AKM estimates with different restrictions on age effects}
\begin{center}
\begin{figure}
\includegraphics[scale=1]{Cardetal/figures/tab3-1.pdf}
\end{figure} 
\end{center}
\end{frame}

\newpage
\begin{frame}{AKM estimates with different restrictions on age effects}
\begin{center}
\begin{figure}
\includegraphics[scale=1]{Cardetal/figures/tab3-2.pdf}
\end{figure} 
\end{center}
\end{frame}

%%%%%%%%%%%%%%%%%%%%%%%%%%%%%%
\subsection*{Worker-Firm Sorting and Limited Mobility Bias}
%%%%%%%%%%%%%%%%%%%%%%%%%%%%%%

%%%%%%%%%%%%%%%%%%%%%%%%%%%%%%
\subsection*{Exogenous Mobility}
%%%%%%%%%%%%%%%%%%%%%%%%%%%%%%

\begin{frame}{Limited mobility bias and exogenous mobility}
\begin{itemize}
\item Under the AKM assumptions, FE estimates are unbiased:
    \begin{equation*}
	 \forall j: \mathbb{E} \left[ \hat{\psi_{j}} \right] = \psi_{j} \text{   and   }  \forall i: \mathbb{E} \left[ \hat{\alpha_{i}} \right] = \alpha_{i}
    \end{equation*}
\item But noise leads to upward bias in variance of FE estimated effects and limited mobility to downward bias in covariance \medskip
\item Exogenous mobility:
    \begin{equation*}
	P(J(i,t)=j | \alpha_{i}, \psi, \epsilon_{i,1},...,\epsilon_{i,T})=P(J(i,t)=j | \alpha_{i}, \psi)
    \end{equation*}
which is an identifying assumption.
\end{itemize}
\end{frame}

\newpage
\begin{frame}{Wage changes for job switchers: men}
\begin{center}
\begin{figure}
\includegraphics[scale=0.85]{Cardetal/figures/fig3.pdf}
\end{figure} 
\end{center}
\end{frame}

\newpage
\begin{frame}{Wage changes for job switchers: women}
\begin{center}
\begin{figure}
\includegraphics[scale=0.85]{Cardetal/figures/fig4.pdf}
\end{figure} 
\end{center}
\end{frame}

\newpage
\begin{frame}{Testing symmetry between switches: men}
\begin{center}
\begin{figure}
\includegraphics[scale=0.95]{Cardetal/figures/fig5-A.pdf}
\end{figure} 
\end{center}
\end{frame}

\newpage
\begin{frame}{Testing symmetry between switches: women}
\begin{center}
\begin{figure}
\includegraphics[scale=0.95]{Cardetal/figures/fig5-B.pdf}
\end{figure} 
\end{center}
\end{frame}

%%%%%%%%%%%%%%%%%%%%%%%%%%%%%%
\subsection*{Additive Separability}
%%%%%%%%%%%%%%%%%%%%%%%%%%%%%%

\newpage
\begin{frame}{Testing additive separability: men}
\begin{center}
\begin{figure}
\includegraphics[scale=0.7]{Cardetal/figures/fig6.pdf}
\end{figure} 
\end{center}
\end{frame}

\newpage
\begin{frame}{Testing additive separability: women}
\begin{center}
\begin{figure}
\includegraphics[scale=0.75]{Cardetal/figures/fig7.pdf}
\end{figure} 
\end{center}
\end{frame}

%%%%%%%%%%%%%%%%%%%%%%%%%%%%%%%%%%%%%%%%%%%%%%%%%%%%%%%%%
\section{Reconciling Rent-Sharing Estimates with Results from Studies of Firm Switching}
%%%%%%%%%%%%%%%%%%%%%%%%%%%%%%%%%%%%%%%%%%%%%%%%%%%%%%%%%

\begin{frame}{Reconciling rent-sharing and AKM estimates}
\begin{itemize}
\item Obtain a rent-sharing elasticity $ \gamma_{w} $ from:
\begin{equation*}
    \ln(w_{it}) = \beta_{w} + \gamma_{w} \ln(VA/N)_{j} + \epsilon_{w,it}
\end{equation*}
\item Obtain $\gamma_{\alpha_{i}}$, $\gamma_{\psi_{j}}$ and $\gamma_{X_{it}' \beta}$ by regressing:
\begin{equation*}
    y = \beta_{y} + \gamma_{y} \ln(VA/N)_{j}  + \epsilon_{y,it}
\end{equation*}
with $y = \alpha_{i}, \psi_{j}$ and $\ln(X_{it}' \beta)$ respectively. \medskip
\item We then have:
\begin{equation*}
    \gamma_{w} = \gamma_{\alpha_{i}} + \gamma_{\psi_{j}} + \gamma_{X_{it}' \beta}  
\end{equation*}
with $ \gamma_{\psi_{j}}$ a clean measure of the true rent-sharing elasticity and $\gamma_{\alpha_{i}} $ and $\gamma_{X_{it}' \beta} $ capturing sorting effects.
\end{itemize}
\end{frame}

\newpage
\begin{frame}{Reconciling rent-sharing and AKM estimates}
\begin{center}
\begin{figure}
\includegraphics[scale=0.95]{Cardetal/figures/tab4-1.pdf}
\end{figure} 
\end{center}
\end{frame}

\newpage
\begin{frame}{Reconciling rent-sharing and AKM estimates}
\begin{center}
\begin{figure}
\includegraphics[scale=0.95]{Cardetal/figures/tab4-2.pdf}
\end{figure} 
\end{center}
\end{frame}

%%%%%%%%%%%%%%%%%%%%%%%%%%%%%%
\subsection*{Differential Rent Sharing}
%%%%%%%%%%%%%%%%%%%%%%%%%%%%%%

\newpage
\begin{frame}{No differential rent-sharing}
\begin{center}
\begin{figure}
\includegraphics[scale=0.85]{Cardetal/figures/fig8.pdf}
\end{figure} 
\end{center}
\end{frame}

%%%%%%%%%%%%%%%%%%%%%%%%%%%%%%%%%%%%%%%%%%%%%%%%%%%%%%%%%
\section{Imperfectly Competitive Labor Markets and Inequality}
%%%%%%%%%%%%%%%%%%%%%%%%%%%%%%%%%%%%%%%%%%%%%%%%%%%%%%%%%

%%%%%%%%%%%%%%%%%%%%%%%%%%%%%%
\subsection*{Market Structure}
%%%%%%%%%%%%%%%%%%%%%%%%%%%%%%

\begin{frame}{Worker utility}
\begin{itemize}
\item For worker $i$ in skill group $S \in \{ L,H \}$ utility of working at firm $j$ is given by:
\begin{equation*}
    u_{iSj} = \beta_{S} \ln(w_{j} - b_{s}) + a_{Sj} + \epsilon_{iSj}
\end{equation*}
with $b_{s}$ reference wage, $ a_{Sj}$ firm-specific amenities valued by all workers, $\epsilon_{iSj} $ idiosyncratic preference for working at firm $j$. \medskip
\item $ \{\epsilon_{iSj} \} $ are independent draws from a type I Extreme Value distribution (to model worker's choice for $j$ as in BLP). \medskip
\item Firms post wages knowing the distribution of $ \{\epsilon_{iSj} \} $ but not individual specific values and workers choose a $j$. 
\end{itemize}
\end{frame}

\begin{frame}{Worker sorting across firms}
\begin{itemize}
\item Given posted wages, workers have logit choice probabilities:
\begin{equation*}
    p_{Sj} = \frac{\exp(\beta_{S}(\ln(w_{Sj} - b_{S}) + a_{Sj}))}{\sum_{k=1}^{J} \exp(\beta_{S}(\ln(w_{Sk} - b_{S}) + a_{Sk}))}
\end{equation*}
\item Assume that $J$ is very large such that:
\begin{equation*}
    p_{Sj} \approx \lambda_{S} \exp(\beta_{S}(\ln(w_{Sj} - b_{S}) + a_{Sj}))
\end{equation*}
\item This gives firm-level supply functions:
\begin{align*}
    \ln(L_{j}(w_{Lj})) & = \ln(\mathcal{L} \lambda_{L}) + \beta_{L} \ln(w_{Lj} - b_{L}) + a_{Lj} \tag{5} \\
    \ln(H_{j}(w_{Lj})) & = \ln(\mathcal{H} \lambda_{H}) + \beta_{H} \ln(w_{Hj} - b_{H}) + a_{Hj} \tag{6} \\
\end{align*}
\end{itemize}
\end{frame}

\begin{frame}{Elasticities of firm-level labor supply}
\begin{itemize}
\item The elasticities of firm-level labor supply are given by:
\begin{align*}
    e_{Lj} & \equiv \frac{\partial \ln(L_{j}(w_{Lj}))}{\partial \ln(w_{Lj})} = \beta_{L} \frac{w_{Lj}}{w_{Lj} - b_{L}} \\
    e_{Hj} & \equiv \frac{\partial \ln(H_{j}(w_{Hj}))}{\partial \ln(w_{Hj})} = \beta_{H} \frac{w_{Hj}}{w_{Hj} - b_{H}}
\end{align*}
\item Finite $ \beta_{S} $ implies workers' idiosyncratic preferences for amenities matter giving firms monopsony power. \medskip
\item As $ \beta_{L}, \beta_{H} \rightarrow \infty $ firm-level labor supply becomes perfectly elastic and we approach competitive labor market.
\end{itemize}
\end{frame}

%%%%%%%%%%%%%%%%%%%%%%%%%%%%%%
\subsection*{Firm Optimization}
%%%%%%%%%%%%%%%%%%%%%%%%%%%%%%

\begin{frame}{Firm optimization}
\begin{itemize}
\item Firm have production functions:
\begin{equation*}
    Y_{j} = T_{j} f(L_{j},H_{j})
\end{equation*}
\item The firm's optimal wage choices solve for the first-order conditions:
\begin{align*}
    w_{Lj} \frac{1+e_{Lj}}{e_{Lj}} & = T_{j} f_{L} \mu_{j} \tag{8} \\
    w_{Hj} \frac{1+e_{Hj}}{e_{Hj}} & = T_{j} f_{H} \mu_{j} \tag{9}
\end{align*}
with $ e_{Sj} $ the elasticity of labor supply, $ f_{S} $ marginal labor productivity, and $\mu_{j} $ the firm's marginal cost.
\end{itemize}
\end{frame}

\begin{frame}{Wage markdowns from marginal product}
\begin{itemize}
\item Using expressions for firm-level labor supply elasticities gives:
\begin{align*}
    w_{Lj} & = \frac{1}{1 + \beta_{L}}b_{L} + \frac{\beta_{L}}{1 + \beta_{L}} T_{j} f_{L} \mu_{j} \tag{10} \\
    w_{Hj} & = \frac{1}{1 + \beta_{H}}b_{H} + \frac{\beta_{H}}{1 + \beta_{H}} T_{j} f_{H} \mu_{j} \tag{11}
\end{align*}
\item Wage is a weighted average of $b_{S}$ and marginal product. \medskip
\item If $ b_{S} < T_{j} f_{S} \mu_{j} $, wage is marked down from marginal product. \medskip
\item As $ \beta_{S} \rightarrow \infty $ wage converges to marginal product and labor market is perfectly competitive.
\end{itemize}
\end{frame}

\begin{frame}{Example for a single skill group}
\begin{itemize}
\item Firm-level labor supply is given by:
\begin{equation*}
    w_{j}(S_{j}) = S_{j}^{\frac{1}{\beta}}
\end{equation*}
\item Output is given by:
\begin{equation*}
    Y_{j}(S_{j}) = T_{j} S_{j}^{\frac{\epsilon - 1}{\epsilon}}
\end{equation*}
\item Product market is perfectly competitive with $P_{j}=1$. \medskip
\item Firms maximize profits:
\begin{equation*}
    \max_{S_{j}} \Pi = P_{j} Y_{j}(S_{j}) - w_{j}(S_{j}) S_{j} = T_{j}S_{j}^{\frac{\epsilon - 1}{\epsilon}} - S_{j}^{\frac{1 + \beta}{\beta}} 
\end{equation*}
\end{itemize}
\end{frame}

\begin{frame}{Example for a single skill group}
\begin{itemize}
\item First-order condition is given by:
\begin{equation*}
    \underbrace{T_{j}\frac{\epsilon - 1}{\epsilon} S_{j}^{-\frac{1}{\epsilon}}}_{MRP} = \underbrace{\frac{1 + \beta}{\beta} S_{j}^{\frac{1}{\beta}}}_{MFC} > S_{j}^{\frac{1}{\beta}} = w_{j}(S_{j})
\end{equation*}
\item Wage is marked down from MRP:
\begin{equation*}
    w_{j} = \frac{\beta}{1+\beta} MRP
\end{equation*}
\item Equilibrium employment is given by:
\begin{equation*}
    S_{j} = \left[ T_{j} \frac{\epsilon - 1}{\epsilon} \frac{\beta}{1+\beta} \right]^{\frac{\beta \epsilon}{\beta + \epsilon}}
\end{equation*}
\item Equilibrium wage is given by:
\begin{equation*}
    w_{j} = \left[ T_{j} \frac{\epsilon-1}{\epsilon} \frac{\beta}{1+\beta} \right]^{\frac{\epsilon}{\epsilon+\beta}}
\end{equation*}
\end{itemize}
\end{frame}

\newpage
\begin{frame}{Equilibrium and wage pass-through of productivity shocks}
\begin{center}
\begin{figure}
\includegraphics[scale=0.85]{Cardetal/figures/fig9.png}
\end{figure} 
\end{center}
\end{frame}

%%%%%%%%%%%%%%%%%%%%%%%%%%%%%%
\subsection*{Baseline Case: Linear Production Function and Fixed Output Price}
%%%%%%%%%%%%%%%%%%%%%%%%%%%%%%

\begin{frame}{Baseline case: linear production function and fixed output price}
\begin{itemize}
\item Return to setup with 2 skill types $S=L,H$. \medskip
\item Linear production function with distribution parameter $ \theta $: \medskip
\begin{equation*}
    Y_{j} = T_{j} N_{j} = T_{j} ( (1 - \theta) L_{j} + \theta H_{j})
\end{equation*}
\item Fixed output price $ P_{j}^{0}$. \medskip
\item The wage equations are given by:
\begin{align*}
    w_{Lj} & = \frac{1}{1 + \beta_{L}}b_{L} + \frac{\beta_{L}}{1 + \beta_{L}} T_{j} P{j}^{0} (1 - \theta) \\
    w_{Hj} & = \frac{1}{1 + \beta_{H}}b_{H} + \frac{\beta_{H}}{1 + \beta_{H}} T_{j} P{j}^{0} \theta
\end{align*}
\end{itemize}
\end{frame}

\begin{frame}{Re-writing the wage equations to capture rent sharing}
\begin{itemize}
\item An alternative way to write the wage equations is:
\begin{align*}
    w_{Lj} & = \frac{b_{L}}{1 + \beta_{L}} \left[ 1 + \frac{\beta_{L}}{b_{L}} T_{j} P{j}^{0} (1 - \theta) \right] \\
    w_{Hj} & = \frac{b_{H}}{1 + \beta_{H}} \left[ 1 + \frac{\beta_{H}}{b_{H}} T_{j} P{j}^{0} \theta \right]
\end{align*}
\item Assume that $b_{L}=(1-\theta)b$ and $b_{H}=\theta b$ \medskip
\item Taking logs gives:
\begin{align*}
    \ln(w_{Lj}) & = \ln \left( \frac{(1-\theta)b}{1 + \beta_{L}} \right) + \ln (1+\beta_{L}R_{j}) \tag{12} \\
    \ln(w_{Hj}) & = \ln \left( \frac{\theta b}{1 + \beta_{H}} \right) + \ln (1+\beta_{H}R_{j}) \tag{13}
\end{align*}
with $ R_{j} \equiv P_{j}^{0}T_{j}/b$ the gap between MRP and $b$.  
\end{itemize}
\end{frame}

\begin{frame}{Rent-sharing elasticities}
\begin{itemize}
\item Value added per standardized unit of labor is:
\begin{equation*}
    \nu_{j} = P_{j}^{0}Y_{j}/N_{j}=P_{j}^{0}T_{j}=R_{j}b
\end{equation*}
\item Rent-sharing elasticities are then given by:
\begin{equation*}
    \xi_{Sj} \equiv \frac{\partial \ln(w_{Sj})}{\partial \ln(\nu_{j})} = \frac{\beta_{S}R_{j}}{1+\beta_{S}R_{j}}
\end{equation*}
which is the same as equation (2). \medskip
\item This elasticity is around 0.1 such that $\beta_{S}R_{j}$ is around 0.1.
\end{itemize}
\end{frame}

\begin{frame}{Elasticities of firm-level labor supply and relative wages}
\begin{itemize}
\item The firm-level labor supply elasticity for group $S$ is given by:
\begin{equation*}
    e_{Sj} = \frac{\beta_{S}w_{Sj}}{w_{Sj}-b_{S}} = \frac{1+\beta_{S}R_{j}}{R_{j}-1}
\end{equation*}
\item If $\beta_{S}R_{j}=0.1$ and $e_{Sj}=4$, we have $R_{j}=1.3$ and $\beta_{S}=0.08$ \medskip
\item $\ln(1+\beta_{S}R_{j}) \approx \beta_{S}R_{j}$ implies that:
\begin{align*}
    \ln \left( \frac{w_{Lj}}{w_{Hj}} \right) = & \ln \left( \frac{\theta}{1-\theta} \right) + \ln \left( \frac{1+\beta_{L}}{1+\beta_{H}} \right) \tag{14}  \\
    & + (\beta_{H} - \beta_{L})R_{j} 
\end{align*}
such that $\beta_{L}=\beta_{H}$ implies that relative wages do not depend on firm-specific productivity shocks.
\end{itemize}
\end{frame}

\begin{frame}{AKM worker and firm effects}
\begin{itemize}
\item We had that:
\begin{align*}
    \ln(w_{Lj}) & \approx \ln \left( \frac{(1-\theta)b}{1 + \beta_{L}} \right) + \beta_{L}R_{j}  \\
    \ln(w_{Hj}) & \approx \ln \left( \frac{\theta b}{1 + \beta_{H}} \right) + \beta_{H}R_{j} 
\end{align*}
\item If $\beta_{L}=\beta_{H}=\beta$, wages can be written as an AKM model:
\begin{equation*}
    \ln(w_{Sj}) = \alpha_{S} + \beta R_{j} = \alpha_{S} + \psi_{j} \tag{15}
\end{equation*}
with $\alpha_{S}$ and $\psi_{j}$ additive worker and firm effects with $ \alpha_{S} \equiv \ln(b/(1+\beta)) + \mathbb{1}(S=L) \ln(1-\theta) + \mathbb{1}(S=H) \ln(\theta)$
\end{itemize}
\end{frame}


\begin{frame}{AKM worker and firm effects}
\begin{itemize}
\item If $\beta_{H} > \beta_{L}$, the log wage gap between $H$ and $L$ will be higher at more profitable firms. \medskip
\item We get AKM model with skill group–specific firm effects:
\begin{equation*}
    \psi_{j}^{H} = \beta_{H}R_{j} > \psi_{j}^{L} = \beta_{L}R_{j} 
\end{equation*}
\item Estimate of $\beta_{H}/ \beta_{L}$ from projecting $\psi_{j}^{H}$ onto $\psi_{j}^{L}$:
\begin{equation*}
    \psi_{j}^{H} = (\beta_{H}/ \beta_{L}) \psi_{j}^{L}
\end{equation*}
\item Card et al. (2016) estimate $\beta$ about 10\% larger for males than females in Portugal (because men value wages more and pass-through of productivity to wages is larger for men).
\end{itemize}
\end{frame}

\begin{frame}{Between-firm sorting}
\begin{itemize}
\item If $\beta_{H}=\beta_{L}$, the market-wide average wage for $S$ is:
\begin{equation*}
    \mathbb{E}[\ln(w_{Si})] = \alpha_{S} + \sum_{j}\psi_{j} \pi_{Sj}
\end{equation*}
with $\pi_{Sj}$ the share of $S$-type workers employed at $j$. \medskip
\item The market-wide wage differential between groups is given by:
\begin{equation*}
    \mathbb{E}[\ln(w_{Hi})] - \mathbb{E}[\ln(w_{Li})] = \alpha_{H} - \alpha_{L} + \sum_{j}\psi_{j} [\pi_{Hj} - \pi_{Lj}] \tag{16}
\end{equation*}
\item Card et al. (2016) show that 15\% of wage differential between men and women in Portugal is explained by men sorting more into firms that pay higher premium (to men and women).
\end{itemize}
\end{frame}

\newpage
\begin{frame}{Between-firm sorting by age and education}
\begin{center}
\begin{figure}
\includegraphics[scale=0.85]{Cardetal/figures/fig10.pdf}
\end{figure} 
\end{center}
\end{frame}

\begin{frame}{Market-wide wage differences when $\beta_{H} \neq \beta_{L}$}
\begin{itemize}
\item If $\beta_{H} \neq \beta_{L}$, the market-wide wage differential is given by:
\begin{align*}
    \mathbb{E}[\ln(w_{Hi})] - \mathbb{E}[\ln(w_{Li})] & = \alpha_{H} - \alpha_{L} \\
    & + \sum_{j}\psi_{j}^{L} [\pi_{Hj} - \pi_{Lj}] + \sum_{j} [\psi_{j}^{H} - \psi_{j}^{L} ]\pi_{Hj}  \\
    &  = \alpha_{H} - \alpha_{L} \\
    & + \sum_{j}\psi_{j}^{H} [\pi_{Hj} - \pi_{Lj}] + \sum_{j} [\psi_{j}^{H} - \psi_{j}^{L} ]\pi_{Lj}
\end{align*}
\item Market-wide wage differential now also accounts for differential rent-sharing within firms. \medskip
\item As in Oaxaca decomposition, there are alternative ways to evaluate the relative importance of between and within firm components.
\end{itemize}
\end{frame}

\begin{frame}{Extensions}
\begin{itemize}
\item Downward sloping firm-specific product demand:
\begin{equation*}
    P_{j} = P_{j}^{0} Y_{j}^{-1/ \epsilon}
\end{equation*}
with $R'_{j} = [(\epsilon - 1)/ \epsilon] \nu_{j}/b$ \medskip
\item Imperfect substitution between skill groups:
\begin{equation*}
    Y_{j} = T_{j} N_{j} = T_{j} \left[ ( (1 - \theta) L_{j}^{\rho} + \theta H_{j}^{\rho}) \right]^{1/\rho}
\end{equation*}
and basic intuition of benchmark model remains \medskip
\item Other extensions would be (non-random) sorting of workers or strategic interactions between employers.
\end{itemize}
\end{frame}

%%%%%%%%%%%%%%%%%%%%%%%%%%%%%%%%%%%%%%%%%%%%%%%%%%%%%%%%%
\section{Conclusions}
%%%%%%%%%%%%%%%%%%%%%%%%%%%%%%%%%%%%%%%%%%%%%%%%%%%%%%%%%

\begin{frame}{Conclusions}
\begin{itemize}
\item CCHK framework in this paper is known as ``new classical monopsony" \medskip
\item It is consistent with evidence on rent-sharing elasticities and AKM models of additive worker and firm effects \medskip
\item New classical monopsony differs from ``modern monopsony" based on search frictions (Burdett\&Mortensen 98) \medskip
\item Interesting extensions would be to bridge gap between both approaches by incorporating dynamics in CCHK. \medskip
\item Other interesting extension would be to incorporate strategic interaction between firms.
\end{itemize}
\end{frame}

\end{document}
